\begin{lemma}
Let \(U,V\) be finite sets and let \(f\) be a map from the ambient vertex set to a label set.  Put \(U_0=U\setminus V\).  Suppose \(\eta:U_0\to \beta\) agrees with \(f\) on \(U_0\).  If, for a real threshold \(\tau\),
\[
  |f(U)\cap f(V)|>|U\cap V|+\tau,
\]
then
\[
  \left\lceil \tau\right\rceil
  \le
  \left|\{v\in V\setminus U: f(v)\in \eta(U_0)\}\right|.
\]
\end{lemma}
\begin{proof}
By \(\noderef{OppositeProjectionIntersectionContainment}\),
\[
  |f(U)\cap f(V)|
  \le
  |U\cap V|+
  \left|\{v\in V\setminus U:f(v)\in f(U\setminus V)\}\right|.
\]
Since \(\eta\) agrees with \(f\) on \(U_0=U\setminus V\), every value in \(f(U\setminus V)\) lies in the image \(\eta(U_0)\).  Hence the displayed exceptional set is contained in
\[
  \{v\in V\setminus U:f(v)\in \eta(U_0)\}.
\]
Combining this containment with the strict lower bound
\(|f(U)\cap f(V)|>|U\cap V|+\tau\) gives
\[
  \tau
  \le
  \left|\{v\in V\setminus U:f(v)\in \eta(U_0)\}\right|.
\]
The defining property of the natural ceiling then yields the claimed
\(\lceil\tau\rceil\) bound.
\end{proof}
